\exercisegroup

\begin{enumerate}
\def\labelenumi{\arabic{enumi}.}
\item
  证明：在集合E上的序关系的集合中，关于“Γ比Γ'更粗”这一序关系而言，极小元恰为E上的全序关系；并且若Γ是E上的任意序关系，则Γ的图是E上所有比Γ更粗的全序关系的图的交集（应用第4节定理2）。由此推出：每个有序集都同构于某个全序集之积的一个子集。
\item
  设 E 为一个有序集，并设 B 为 E 中由诱导序良序的子集所构成的集合。证明 B 上的关系“X 是 Y 的一个段”是 X 与 Y 之间的一个序关系，并且 B 关于该序关系是归纳的。由此推出，E 中存在没有严格上界的良序子集。
\item
  设 E 是一个有序集。证明存在 E 的两个子集 A, B，使得 \(A \cup B = E\) 并且 \(A \cap B = \emptyset\) ，并且使得 A 是良序的而 B 没有最小元素（例如，取 B 为 E 的那些没有最小元素的子集的并集）。* 给出一个例子，其中存在将 E 分成具有这些性质的两个子集的多种分法。*
\item
  一个有序集 F 被称为部分良序的，如果 F 的每个全序子集都是良序的。证明：在每个有序集 E 中都存在一个在 E 中共尾的部分良序子集。（考虑 E 的部分良序子集构成的集 F，以及 F 的元素 X 和 Y 之间的序关系“X ⊂ Y 且 Y --- X 中没有元素被 X 中的任何元素所界定”。证明 F 关于此序关系是归纳的。）
\item
  设 E 是一个有序集，并设 J 是 E 的自由子集的集合，按第 1 节习题 5 中定义的关系排序。证明：若 E 是归纳的，则 J 有最大元。
\end{enumerate}

¶ 6. 设 E 是一个有序集，并设 f 是从 E 到 E 的一个映射，使得 \(f(x) \geqslant x\) 对于所有 \(x \in E\) .

\begin{enumerate}
\def\labelenumi{(\alph{enumi})}
\item
  让 \(\mathfrak{G}\) 为具有以下性质的子集 \(\mathbf{M}\) （\(\mathbf{E}\) 的子集）所构成的集合：（1）关系 \(x \in \mathbf{M}\) 蕴含 \(f(x) \in \mathbf{M}\) ； (2) 若 \(\mathbf{M}\) 的一个非空子集在 \(\mathbf{E}\) 中有一个最小上界，则此最小上界属于 \(\mathbf{M}\) 。对于每一个 \(a \in \mathbf{E}\) ，证明交集 \(\mathbf{C}_a\) （\(\mathfrak{G}\) 中包含 \(a\) 的集合的交集）也属于 \(\mathfrak{G}\) ；即 \(\mathbf{C}_a\) 是良序的；并且若 \(\mathbf{C}_a\) 有一个上界 \(b\) （在 \(\mathbf{E}\) 中），那么 \(b \in \mathbf{C}_a\) 并且 \(f(b) = b\) 。 \(\mathbf{C}_a\) 被称为 \(a\) 的链（关于函数 \(f\) ）。（考虑集合 \(\mathfrak{M}\) ，其元素为空集和满足以下条件的子集 \(\mathbf{X}\) （\(\mathbf{E}\) 的子集）：包含 \(a\) 并且有一个最小上界 \(m\) （在 \(\mathbf{E}\) 中），使得要么 \(m \notin \mathbf{X}\) 要么 \(f(m) > m\) ；将第3节中的引理3应用于集合 \(\mathfrak{M}\) 。）
\item
  从（a）推出，如果 \(\mathbf{E}\) 是归纳的，则存在 \(b \in \mathbf{E}\) 使得 \(f(b) = b\) .
\end{enumerate}

\begin{enumerate}
\def\labelenumi{\arabic{enumi}.}
\setcounter{enumi}{6}
\tightlist
\item
  设 E 是一个有序集，并设 F 为 E 中所有闭包（§1，习题13）的集合。对 F 按如下方式排序： \(u \leqslant v\) 每当 \(u(x) \leqslant v(x)\) 对于所有 \(x \in E\) 成立。则 F 有最小元 \(e\) ，即 E 到自身的恒等映射。对每个 \(u \in F\) ，设 I(u) 表示 E 中在 \(u\) 下不变的元素的集合。
\end{enumerate}

\begin{enumerate}
\def\labelenumi{(\alph{enumi})}
\item
  证明 \(u \leqslant v\) （在 \(\mathbf{F}\) 中）当且仅当 \(\mathbf{I}(v) \subset \mathbf{I}(u)\) 。
\item
  证明：如果E的每一对元素在E中都有一个最大下界，那么F的每一对元素在F中也有一个最大下界。若E是完全格，则F也是完全格（§1，习题11）。
\item
  证明若 E 是归纳的（关于关系 \(\leqslant\) ），则 F 中任意一对元素 u, v 在 F 中都有最小上界。（证明若 \(f(x) = v(u(x))\) 并且 \(w(x)\) 表示x相对于f的链的最大元素（练习6），则w是E中的一个闭包，并且是u和v的最小上界。）
\end{enumerate}

¶ 8. 一个有序集 E 被称为（在右侧）分歧的，如果对于 E 中满足 x \textless{} y 的每一对元素 x, y，存在 z \textgreater{} x，使得 y 与 z 不可比较。若 E 是分歧的（在右侧）且没有极大元，则称 E 是完全分歧的（在右侧）。每个反定向集（§1，习题22）都是分歧的。

\begin{enumerate}
\def\labelenumi{(\alph{enumi})}
\item
  设E是一个有序集，并设a是E的一个元素。设 \(R_{a}\) 表示E中以a为最小元素的分歧子集的集合。证明 \(R_{a}\) 在包含序下具有极大元素。
\item
  如果 \(\mathbf{E}\) 是分支的（§1，习题24），证明 \(\mathfrak{R}_a\) 的每个极大元素是完全分歧的。
\item
  给出一个分支但非分歧的集合的例子。§1习题24(c)中定义的分支集是完全分歧的。
\item
  设E是一个集合，其中每个区间 \(\leftarrow, x]\) 是全序的。证明 E 有一个反定向的共尾子集（§1，练习 22）（利用 (b)）。
\end{enumerate}

\begin{enumerate}
\def\labelenumi{\arabic{enumi}.}
\setcounter{enumi}{8}
\item
  序数和 \(\sum_{i\in I}E_i\) ( \(\S 1\) ，练习 3）是良序的当且仅当 \(I\) 与每个 \(E_i\) 都是良序的。
\item
  令 I 为一个有序集，且令 \((\mathbf{E}_t)_{t \in I}\) 为一个有序集族，其中所有集合都等于同一个有序集E。证明序数和 \(\sum_{t \in I} \mathbf{E}_t\) 若 \(t \in I\) ，则序集由这些序的和给出 \(\mathbf{E}_t\) （§1，习题3）与序列 \((\mathbf{F}_{\lambda})_{\lambda \in |\alpha, \beta|}\) 的字典序乘积同构，其中集合 \(\{\alpha, \beta\}\) （含两个不同元素）由图形为 \(\{(\alpha, \alpha), (\alpha, \beta), (\beta, \beta)\}\) 的关系良序化，并且其中 \(\mathbf{F}_{\alpha} = \mathbf{I}\) 并且 \(\mathbf{F}_{\beta} = \mathbf{E}\) 。该乘积称为 \(\mathbf{E}\) 由 \(\mathbf{I}\) 的字典序乘积，记作 E.I.
\end{enumerate}

¶ * 11. 设 I 是一个良序集，并设 \((E_t)_{t \in I}\) 是一族有序集，其中每个有序集至少包含两个不同的可比较元素。则 \(E_t\) 的字典序乘积是良序的当且仅当每个 \(E_t\) 都是良序的且 I 是有限的（若 I 是无限的，则在 \(E_t\) 的字典序乘积中构造一个严格递减的无限序列）。*

\begin{enumerate}
\def\labelenumi{\arabic{enumi}.}
\setcounter{enumi}{11}
\tightlist
\item
  设 I 是一个全序集，并设 \((\mathbf{E}_{t})_{t\in\mathbf{I}}\) 为一个由I指标化的有序集族。令 \(R\{x,y\}\) 表示 \(E=\prod_{i\in I}E_{i}\) 上的以下关系：“指标 \(i\in I\) （使得 \(pr_{i}x\neq pr_{i}y\) ）所成的集合是良
\end{enumerate}

有序的，并且如果 \(x\) 是 I 的这个子集的最小元素，则我们有 \(\mathrm{pr}_x x < \mathrm{pr}_x y\) 。证明 \(R\{x, y\}\) 是 E 上 \(x\) 与 \(y\) 之间的一个序关系。如果这些 \(E_t\) 是全序的，证明 E 关于关系“ \(x\) 和 \(y\) 是可比较的”（第二章，§6，习题10）的连通分量是全序集。假设每个 \(E_t\) 至少有两个元素。那么E是全序的当且仅当I是良序的且每个 \(E_t\) 是全序的（使用习题3）；并且E此时是诸 \(E_t\) 的字典序乘积。

\begin{enumerate}
\def\labelenumi{\arabic{enumi}.}
\setcounter{enumi}{12}
\tightlist
\item
  \begin{enumerate}
  \def\labelenumii{(\alph{enumii})}
  \tightlist
  \item
    让 \(\operatorname{Is}(\Gamma, \Gamma')\) 为关系“ \(\Gamma\) 是（E 上的）一个序，且 \(\Gamma'\) 是（E' 上的）一个序，且存在一个从按 \(\Gamma\) 排序的 E 到按 \(\Gamma'\) 排序的 E' 上的一个同构”。证明 \(\operatorname{Is}(\Gamma, \Gamma')\) 是每个元素为序关系的集合上的一个等价关系。术语 \(\tau_{\Delta}(\operatorname{Is}(\Gamma, \Delta))\) 是一个序，称为 \(\Gamma\) 的序型，并且用 \(\operatorname{Ord}(\Gamma)\) 表示，或 \(\operatorname{Ord}(\mathbf{E})\) 通过滥用记号。两个有序集同构当且仅当它们的序型相等。
  \end{enumerate}
\end{enumerate}

\begin{enumerate}
\def\labelenumi{(\alph{enumi})}
\setcounter{enumi}{1}
\item
  让 \(\mathbb{R}\{\lambda, \mu\}\) 是如下关系：“ \(\lambda\) 是一个序型，且 \(\mu\) 是一个序型，且存在一个从按 \(\lambda\) 排序的集合到按 \(\mu\) 排序的集合的一个子集上的同构”。证明 \(\mathbb{R}\{\lambda, \mu\}\) 是 \(\lambda\) 与 \(\mu\) 之间的预序关系。它将被记为 \(\lambda \prec \mu\) 。
\item
  令 I 为一个有序集，且令 \((\lambda_{t})_{t\in I}\) 为一个由I指标化的序型族。由按 \(\lambda_{t}(t\in I)\) 排序的集合所成的族（§1，习题3）的序数和的序型被称为诸序型 \(\lambda_{t}(t\in I)\) 的序数和，并且用 \(\sum_{i\in I}\lambda_{i}\) 表示。如果 \((E_{i})_{i\in I}\) 是一个有序集的族，则 \(\sum_{i\in I}E_{i}\) 的序型是 \(\sum_{i\in I}\operatorname{Ord}(E_{i})\) 。如果 I 是一个族 \((J_{x})_{x\in K}\) 的序数和，证明
\end{enumerate}

\[
\sum_ {x \in K} \left(\sum_ {i \in J _ {x}} \lambda_ {i}\right) = \sum_ {i \in I} \lambda_ {i}.
\]

\begin{enumerate}
\def\labelenumi{(\alph{enumi})}
\setcounter{enumi}{3}
\tightlist
\item
  设 I 是一个良序集，且 \((\lambda_{t})_{i\in I}\) 是一个由I索引的序型族。由按 \(\lambda_{t} (t\in I)\) 排序的集合所成的族的字典序乘积的序型被称为该族 \((\lambda_{t})_{i\in I}\) 的序数积，并且用 \(\mathsf{P}_{i\in I}\lambda_{t}\) 表示。如果 \((E_{t})_{i\in I}\) 是一个有序集族，族 \((E_{t})_{i\in I}\) 的字典序乘积的序型是 \(\mathsf{P}_{i\in I}\operatorname{Ord}(E_{t})\) 。如果 I 是一族良序集 \((J_{x})_{x\in K}\) （由一个良序集 \(K\) 索引）的序数和，证明
\end{enumerate}

\[
\underset {x \in \mathbf {K}} {\mathsf {P}} \left(\underset {i \in J _ {x}} {\mathsf {P}} \lambda_ {i}\right) = \underset {i \in I} {\mathsf {P}} \lambda_ {i}.
\]

\begin{enumerate}
\def\labelenumi{(\alph{enumi})}
\setcounter{enumi}{4}
\item
  我们用 \(\lambda + \mu\) （相应地， \(\mu\lambda\) ）表示族 \((\xi_{i})_{i \in J}\) 的序数和（相应地，序数积），其中 \(J = \{\alpha, \beta\}\) 是一个具有两个不同元素的集合，其序关系的图是 \(\{(\alpha, \alpha), (\alpha, \beta), (\beta, \beta)\}\) ，并且其中 \(\xi_{\alpha} = \lambda\) , \(\xi_{\beta} = \mu\) 。证明：若 I 是一个序型为 \(\lambda\) 的良序集，并且如果 \((\mu_{i})_{i \in I}\) 是序型的一个族，使得 \(\mu_{i} = \mu\) 对于每一个 \(i \in I\) ，那么 \(\sum_{i \in I} \mu_{i} = \mu\lambda\) 。我们有 \((\lambda + \mu) + \nu = \lambda + (\mu + \nu)\) , \((\lambda\mu)\nu = \lambda(\mu\nu)\) ，以及 \(\lambda(\mu + \nu) = \lambda\mu + \lambda\nu\) （但一般地， \(\lambda + \mu \neq \mu + \lambda\) , \(\lambda\mu \neq \mu\lambda\) , \((\lambda + \mu)\nu \neq \lambda\nu + \mu\nu\) ).
\item
  让 \((\lambda_{t})_{t\in I}\) 和 \((\mu_{t})_{t\in I}\) 是由同一有序集 \(I\) 索引的两个序型族。证明：如果 \(\lambda_{t} \prec \mu_{t}\) 对于每一个 \(t \in I\) ，那么 \(\sum_{i \in I} \lambda_{t} \prec \sum_{i \in I} \mu_{t}\) 并且（如果 \(I\) 是良序的） \(\mathsf{P}_{t} \lambda_{t} \prec \mathsf{P}_{t \in I} \mu_{t}\) 。如果 \(J\) 是 \(I\) 的一个子集，证明 \(\sum_{i \in J} \lambda_{t} \prec \sum_{i \in I} \lambda_{t}\) 并且（如果 \(I\) 是良序的，且 \(\lambda_{t}\) 是非空的） \(\mathsf{P}_{t \in J} \lambda_{t} \prec \mathsf{P}_{t \in I} \lambda_{t}\) 。
\item
  让 \(\lambda^{*}\) 表示按与序 \(\lambda\) 相反的序关系排序的集合的序型。于是我们有
\end{enumerate}

\[
(\lambda^ {*}) ^ {*} = \lambda \quad \text { 且 } \quad \left(\sum_ {i \in I} \lambda_ {i}\right) ^ {*} = \sum_ {i \in I ^ {*}} \lambda_ {i},
\]

其中 I* 表示集合 I 赋予与 I 上给定序关系相反的序关系。

\begin{enumerate}
\def\labelenumi{\arabic{enumi}.}
\setcounter{enumi}{13}
\tightlist
\item
  序数是一个良序集合的序型（习题13）。
\end{enumerate}

\begin{enumerate}
\def\labelenumi{(\alph{enumi})}
\tightlist
\item
  证明：如果 \((\lambda_{i})_{i\in I}\) 是由良序集合索引的序数族 \(I\) ，则序数和 \(\sum_{i\in I}\lambda_{i}\) 是一个序数；* 且此外，如果 \(I\) 是有限的，则序数乘积 \(\mathsf{P}_{i\in I}\lambda_{i}\) 是一个序数（习题11）。* 空集的序型记为0，而含一个元素的集合的序型记为1（通过滥用语言，参见§3）。证明
\end{enumerate}

\[
\alpha + 0 = 0 + \alpha = \alpha \quad \text { 且 } \quad \alpha . 1 = 1. \alpha = \alpha
\]

对于每一个序数 \(\alpha\) .

\begin{enumerate}
\def\labelenumi{(\alph{enumi})}
\setcounter{enumi}{1}
\item
  证明关系“ \(\lambda\) 是序数且 \(\mu\) 是序数且 \(\lambda \prec \mu\) '' 是一个良序关系，记作 \(\lambda \leqslant \mu\) 。 （注意，如果 \(\lambda\) 和 \(\mu\) 是序数，则关系 \(\lambda \prec \mu\) 等价于“ \(\lambda\) 等于 \(\mu\) 的某个段（segment）的序型”（见第5号，定理3，推论3）：给定一个序数的族 \((\lambda_{i})_{i \in I}\) ，考虑一个建立在 \(I\) 上的良序，并取由 \(\lambda_{i}\) 排序的集合族的序数和；最后，使用第1节中的命题2。）
\item
  设 \(\alpha\) 是一个序数。证明关系“ \(\xi\) 是序数且 \(\xi \leqslant \alpha\) ”关于 \(\xi\) 是收集性的，并且集合 \(O_{\alpha}\) （所有小于 \(\alpha\) 的序数所成）是良序集，满足 \(\operatorname{Ord}(O_{\alpha}) = \alpha\) 。以后常把 \(O_{\alpha}\) 与 \(\alpha\) 等同起来。
\item
  证明：对于每一个序数族 \((\xi_{i})_{i\in I}\) 存在唯一一个序数 \(\alpha\) 使得关系“ \(\lambda\) 是序数且 \(\xi_{i}\leqslant\lambda\) 对于所有 \(i\in I\) ”等价于 \(\alpha\leqslant\lambda\) 。通过滥用语言， \(\alpha\) 被称为序数族 \((\xi_{i})_{i\in I}\) 的最小上界，并记作 \(\alpha=\sup_{i\in I}\xi_{i}\)
\end{enumerate}

（它是 \(\{\alpha\}\) 与诸 \(\xi_{i}\) 所成集合的并集的最大元素）。诸序数 \(\xi < \alpha\) 所成集合的最小上界要么是 \(\alpha\) 要么是一个序数 \(\beta\) （使得 \(\alpha = \beta + 1\) ）；在后一种情况下， \(\beta\) 被称为 \(\alpha\) 的前驱。

\begin{enumerate}
\def\labelenumi{\arabic{enumi}.}
\setcounter{enumi}{14}
\tightlist
\item
  \begin{enumerate}
  \def\labelenumii{(\alph{enumii})}
  \tightlist
  \item
    让 \(\alpha\) 和 \(\beta\) 是两个序数。证明不等式 \(\alpha < \beta\) 等价于 \(\alpha + 1 \leqslant \beta\) ，并且它蕴含不等式 \(\xi + \alpha < \xi + \beta\) , \(\alpha + \xi \leqslant \beta + \xi\) , \(\alpha\xi \leqslant \beta\xi\) 对于所有序数 \(\xi\) ，以及 \(\xi\alpha < \xi\beta\) （若 \(\xi > 0\) ）。
  \end{enumerate}
\end{enumerate}

\begin{enumerate}
\def\labelenumi{(\alph{enumi})}
\setcounter{enumi}{1}
\item
  从（a）推出不存在包含所有序数的集合（利用习题14（d））。
\item
  让 \(\alpha, \beta, \mu\) 是三个序数。证明每一个关系 \(\mu + \alpha < \mu + \beta\) , \(\alpha + \mu < \beta + \mu\) 蕴含 \(\alpha < \beta\) ；并且每一个关系 \(\mu \alpha < \mu \beta\) , \(\alpha \mu < \beta \mu\) 蕴含 \(\alpha < \beta\) 只要 \(\mu > 0\) .
\item
  证明关系 \(\mu + \alpha = \mu + \beta\) 蕴含 \(\alpha = \beta\) ，以及 \(\mu \alpha = \mu \beta\) 蕴含 \(\alpha = \beta\) 只要 \(\mu > 0\) .
\item
  两个序数 \(\alpha\) 和 \(\beta\) 满足 \(\alpha \leqslant \beta\) 当且仅当存在一个序数 \(\xi\) 使得 \(\beta = \alpha + \xi\) 。这个序数 \(\xi\) 是唯一的，并且满足 \(\xi \leqslant \beta\) ；它记作 \((- \alpha) + \beta\) 。
\item
  让 \(\alpha, \beta, \zeta\) 是三个序数，使得 \(\zeta < \alpha\beta\) 。证明存在两个序数 \(\xi, \eta\) 使得 \(\zeta = \alpha\eta + \xi\) 并且 \(\xi < \alpha, \eta < \beta\) （参见第5号，定理3，推论3）。此外， \(\xi\) 和 \(\eta\) 由这些条件唯一确定。
\end{enumerate}

\begin{enumerate}
\def\labelenumi{\arabic{enumi}.}
\setcounter{enumi}{15}
\tightlist
\item
  一个序数 \(\rho > 0\) 称为不可分解的，如果不存在一对序数 \(\xi, \eta\) 使得 \(\xi < \rho, \eta < \rho\) ，以及 \(\xi + \eta = \rho\) .
\end{enumerate}

\begin{enumerate}
\def\labelenumi{(\alph{enumi})}
\item
  一个序数 \(\rho\) 是不可分解的，当且仅当 \(\xi + \rho = \rho\) 对于每一个序数 \(\xi\) 使得 \(\xi < \rho\) .
\item
  如果 \(\rho > 1\) 是一个不可分解的序数，并且如果 \(\alpha\) 是任意序数 \(>0\) ，那么 \(\alpha\rho\) 是不可分解的，反之亦然（使用习题15（f））。
\item
  如果 \(\rho\) 是不可分解的，并且如果 \(0 < \alpha < \rho\) ，那么 \(\rho = \alpha\xi\) ，其中 \(\xi\) 是一个不可分解的序数（使用习题15（f））。
\item
  让 \(\alpha\) 为一个序数 \(>0\) 。证明在不可分解序数中存在一个最大的不可分解序数 \(\leqslant \alpha\) （考虑分解 \(\alpha = \rho + \xi\) ，其中 \(\rho\) 是不可分解的）。
\item
  如果 \(\mathbf{E}\) 是任意一组不可分解序数，从（d）推出 \(\mathbf{E}\) 的上确界（习题14（d））是一个不可分解的序数。
\end{enumerate}

\begin{enumerate}
\def\labelenumi{\arabic{enumi}.}
\setcounter{enumi}{16}
\tightlist
\item
  给定一个序数 \(\alpha_0\) ，一个项 \(f(\xi)\) 被称为（关于 \(\xi\) 的）序数函数符号，定义于 \(\xi \geqslant \alpha_0\) ，如果关系“ \(\xi\) 是序数且 \(\xi \geqslant \alpha_0\) ”蕴含关系“ \(f(\xi)\) 是一个序数”； \(f(\xi)\) 被称为正规的，如果关系 \(\alpha_0 \leqslant \xi < \eta\) 蕴含 \(f(\xi) < f(\eta)\) 并且，如果对于每个族 \((\xi_t)_{t \in I}\) （由序数 \(\geqslant \alpha_0\) 组成），我们有 \(\sup_{t \in I} f(\xi_t) = f(\sup_{t \in I} \xi_t)\) （参见练习 14 (d)）。
\end{enumerate}

\begin{enumerate}
\def\labelenumi{(\alph{enumi})}
\item
  证明：对于每个序数 \(\alpha > 0\) , \(\alpha + \xi\) 和 \(\alpha\xi\) 是为 \(\xi \geqslant 0\) 定义的正规函数符号（利用练习 15 (f)）。
\item
  让 \(w(\xi)\) 是一个定义在 \(\xi \geqslant \alpha_0\) 上的序数函数符号，使得 \(w(\xi) \geqslant \xi\) 并且使得 \(\alpha_0 \leqslant \xi < \eta\) 蕴含 \(w(\xi) < w(\eta)\) 。又设 \(g(\xi, \eta)\) 是一个项，使得关系“ \(\xi\) 和 \(\eta\) 是序数 \(\geqslant \alpha_0\) ”蕴含关系“ \(g(\xi, \eta)\) 是一个序数，使得 \(g(\xi, \eta) > \xi\) ''.
\end{enumerate}

定义一个项 \(f(\xi, \eta)\) ，具有以下性质：(1) 对每个序数 \(\xi \geqslant \alpha_0\) , \(f(\xi, 1) = w(\xi)\) ；(2) 对每个序数 \(\xi \geqslant \alpha_0\) 以及每个序数 \(\eta > 1\) , \(f(\xi, \eta) = \sup_{0 < \zeta < \eta} g(f(\xi, \zeta), \xi)\) （使用第2节中的准则C60）。证明若 \(f_1(\xi, \eta)\) 是另一个具有这两个性质的项，则 \(f(\xi, \eta) = f_1(\xi, \eta)\) 对于所有 \(\xi \geqslant \alpha_0\) 以及所有 \(\eta \geqslant 1\) 。证明对每个序数 \(\xi \geqslant \alpha_0\) , \(f(\xi, \eta)\) 是关于 \(\eta\) 的正规函数符号（对所有 \(\eta \geqslant 1\) 有定义）。证明 \(f(\xi, \eta) \geqslant \xi\) 对于所有 \(\eta \geqslant 1\) 并且 \(\xi \geqslant \alpha_0\) ，以及 \(f(\xi, \eta) \geqslant \eta\) 对于所有 \(\xi \geqslant \sup (\alpha_0, 1)\) 和 \(\eta \geqslant 1\) 。此外，对每一对序数 \((\alpha, \beta)\) ，使得 \(\alpha > 0\) , \(\alpha \geqslant \alpha_0\) ，以及 \(\beta \geqslant w(\alpha)\) ，存在唯一一个序数 \(\xi\) 使得

\[
f (\alpha , \xi) \leqslant \beta <   f (\alpha , \xi + 1),
\]

并且我们有 \(\xi \leqslant \beta\) .

\begin{enumerate}
\def\labelenumi{(\alph{enumi})}
\setcounter{enumi}{2}
\item
  如果我们取 \(\alpha_0 = 0\) , \(w(\xi) = \xi + 1\) , \(g(\xi, \eta) = \xi + 1\) ，那么 \(f(\xi, \eta) = \xi + \eta\) 。如果我们取 \(\alpha_0 = 1\) , \(w(\xi) = \xi\) , \(g(\xi, \eta) = \xi + \eta\) ，那么 \(f(\xi, \eta) = \xi \eta\) .
\item
  证明：如果关系 \(\alpha_0 \leqslant \xi \leqslant \xi'\) , \(\alpha_0 \leqslant \eta \leqslant \eta'\) 蕴含 \(g(\xi, \eta) \leqslant g(\xi', \eta')\) ，则关系 \(\alpha_0 \leqslant \xi \leqslant \xi'\) , \(1 \leqslant \eta \leqslant \eta'\) 蕴含 \(f(\xi, \eta) \leqslant f(\xi', \eta')\) 。如果关系 \(\alpha_0 \leqslant \xi \leqslant \xi'\) , \(\alpha_0 \leqslant \eta < \eta'\) 蕴含 \(g(\xi, \eta) < g(\xi, \eta')\) 并且 \(g(\xi, \eta) \leqslant g(\xi', \eta)\) ，则关系 \(\alpha_0 \leqslant \xi < \xi'\) 并且 \(\eta \geqslant 0\) 蕴含 \(f(\xi, \eta + 1) < f(\xi', \eta + 1)\) .
\item
  假设 \(w(\xi) = \xi\) 并且关系 \(\alpha_0 \leqslant \xi \leqslant \xi'\) , \(\alpha_0 \leqslant \eta < \eta'\) 蕴含 \(g(\xi, \eta) < g(\xi, \eta')\) 并且 \(g(\xi, \eta) \leqslant g(\xi', \eta)\) 。进一步假设，对于每个 \(\xi \geqslant \alpha_0\) , \(g(\xi, \eta)\) 是关于 \(\eta\) 的正规函数符号（定义于 \(\eta \geqslant \alpha_0\) ），并且，每当 \(\xi \geqslant \alpha_0\) , \(\eta \geqslant \alpha_0\) ，以及 \(\zeta \geqslant \alpha_0\) 时，我们有结合性关系
\end{enumerate}

\[
g (g (\xi , \eta), \zeta) = g (\xi , g (\eta , \zeta)).
\]

证明：如果 \(\xi \geqslant \alpha_0\) , \(\eta \geqslant 1\) ，以及 \(\zeta \geqslant 1\) ，那么

\[
g (f (\xi , \eta), f (\xi , \zeta)) = f (\xi , \eta + \zeta)
\]

（ \(g\) 关于 \(f\) 的“分配性”）且

\[
f (f (\xi , \eta), \zeta) = f (\xi , \eta \zeta)
\]

（ \(f\) 的“结合性”）。

¶ 18. 在练习17（b）中定义的过程中，取 \(\alpha_{0}=1+1\) （滥用记号，记为2），

\[
w (\xi) = \xi , \quad g (\xi , \eta) = \xi \eta .
\]

把 \(f(\xi, \eta)\) 记作 \(\xi^{\eta}\) ，并定义 \(\alpha^{0}\) 为1（对所有序数 \(\alpha\) ）；同时定义 \(0^{\beta}\) 为0，且 \(1^{\beta}\) 为1（对所有序数 \(\beta \geqslant 1\) ）。

\begin{enumerate}
\def\labelenumi{(\alph{enumi})}
\item
  证明如果 \(\alpha > 1\) 并且 \(\beta < \beta'\) ，我们有 \(\alpha^{\beta} < \alpha^{\beta'}\) 并且，对于每个序数 \(\alpha > 1\) , \(\alpha^{\xi}\) 是关于 \(\xi\) 的正规函数符号。此外，如果 \(0 < \alpha \leqslant \alpha'\) ，我们有 \(\alpha^{\beta} \leqslant \alpha'^{\beta}\) .
\item
  证明 \(\alpha^{\xi}.\alpha^{\eta} = \alpha^{\xi +\eta}\) 并且 \((\alpha^{\xi})^{\eta} = \alpha^{\xi \eta}\) .
\item
  证明：如果 \(\alpha \geqslant 2\) 并且 \(\beta \geqslant 1\) , \(\alpha^{\beta} \geqslant \alpha \beta\) .
\item
  对于每一对序数 \(\beta \geqslant 1\) 并且 \(\alpha \geqslant 2\) ，存在三个序数 \(\xi, \gamma, \delta\) 使得 \(\beta = \alpha^{\xi}\gamma + \delta\) ，其中 \(0 < \gamma < \alpha\) 并且 \(\delta < \alpha^{\xi}\) ，并且这三个序数由这些条件唯一确定。
\end{enumerate}

* 19. 设 \(\alpha\) 和 \(\beta\) 为两个序数，并设 E, F 为两个良序集，使得 \(\operatorname{Ord}(E) = \alpha\) 和 \(\operatorname{Ord}(F) = \beta\) 。在从 F 到 E 的映射所成的集合 \(E^F\) 中，考虑由满足以下条件的映射 \(g\) 组成的子集 G： \(g(y)\) 对于除有限个元素外的所有元素 \(y \in F\) 而言等于E的最小元素。如果 \(F^*\) 为将 F 赋予相反序后所得的序集，证明 G 是关于关系“ \(x\) 和 \(y\) 可比较”（第二章，第6节，习题10）在乘积 \(E^{F*}\) （赋予习题12中定义的序）中的一个连通分量，并证明 G 是良序的。此外，证明 \(\operatorname{Ord}(G) = \alpha^\beta\) （利用习题17(b)的唯一性性质）。

\begin{enumerate}
\def\labelenumi{\arabic{enumi}.}
\setcounter{enumi}{19}
\tightlist
\item
  称集合 X 是传递的，如果关系 \(x \in \mathbf{X}\) 蕴含 \(x \subset \mathbf{X}\) .
\end{enumerate}

\begin{enumerate}
\def\labelenumi{(\alph{enumi})}
\item
  若 Y 是传递集，则 \(Y \cup \{Y\}\) 也是传递集。如果 \((Y_{t})_{t \in I}\) 是传递集族，则 \(\bigcup_{i \in I} Y_{t}\) 和 \(\bigcap_{i \in I} Y_{t}\) 是传递的。
\item
  集合 X 称为伪序数，如果每个满足 \(Y \subset X\) 和 \(Y \neq X\) 的传递集 Y 都是 X 的元素。集合 S 称为体面的，如果关系 \(x \in S\) 蕴含 \(x \notin x\) 。证明每个伪序数都是传递且体面的（考虑 X 的体面传递子集的并，并利用 (a)）。若 X 是伪序数，则 \(X \cup \{X\}\) 也是伪序数。
\item
  设 X 是传递集，并假设每个 \(x \in X\) 是一个伪序数。那么X是一个伪序数（注意，对于每一个 \(x \in X\) , \(x \cup \{x\}\) 是包含在X中的一个伪序数）。
\item
  证明 \(\phi\) 是一个伪序数，并且伪序数X的每个元素都是伪序数。（考虑X的传递子集的并集，这些子集的元素都是伪序数。）
\item
  如果 \((\mathbf{X}_t)_{t\in \mathbf{I}}\) 是伪序数的一个族，则 \(\bigcap_{t\in \mathbf{I}}\mathbf{X}_t\) 是该族的最小元素（关于包含关系）。（利用（b）。）由此推出，若 E 是一个伪序数，则关系 \(x\subset y\) （E 的元素 \(x,y\) 之间的关系）是一个良序关系。
\item
  证明：对于每个序数 \(\alpha\) 存在唯一一个伪序数 \(\mathbf{E}_{\alpha}\) 使得 \(\operatorname{Ord}(\mathbf{E}_{\alpha}) = \alpha\) （利用（e）和准则C60）。特别地，那些序型为 \(0, 1, 2 = 1 + 1\) ，以及 \(3 = 2 + 1\) 分别
\end{enumerate}

\[
\phi , \quad \{\phi \}, \quad \{\phi , \{\phi \} \}, \quad \{\phi , \{\phi \}, \{\phi , \{\phi \} \} \}.
\]

